\section{Limitations connues et recommandations}
\label{sec:limitations}

Certaines limites ont été identifiées mais délibérément documentées plutôt que corrigées dans le cadre de ce projet, soit parce que leur correction dépasse le périmètre d'un ajustement ciblé, soit parce qu'elles relèvent d'un arbitrage produit qui n'a pas encore été tranché.

\textbf{La négation n'est pas comprise par la couche de récupération sémantique.} Une question formulée comme ``quels centres n'acceptent \emph{pas} les enfants de moins de trois ans ?'' obtient une réponse portant sur les centres qui \emph{acceptent} cette tranche d'âge -- l'exact inverse de l'intention. La recherche vectorielle et le classement métier actuels ne modélisent pas la négation logique~; ce cas a été observé et reproduit dans le jeu de quinze questions de test, et reste documenté comme limitation connue plutôt que corrigé, une solution robuste nécessitant une étape de compréhension de la question en amont de la récupération.

\textbf{Le calcul arithmétique confié au modèle de langage sur de grands nombres n'est pas fiable.} Le calculateur métier couvre les opérations simples détectées explicitement, mais toute opération numérique laissée à la charge du modèle de langage lui-même (par exemple une somme demandée implicitement au sein d'une réponse plus large) reste sujette à erreur, comme tout grand modèle de langage généraliste. La recommandation est de systématiser le calcul en code pour toute opération numérique identifiable, plutôt que de faire confiance au modèle.

\textbf{La latence reste la contrainte dominante de l'expérience utilisateur.} Sur ce serveur CPU partagé, une réponse ancrée sur le contexte prend de 60 à plus de 300 secondes selon la contention, avec une variabilité significative observée en conditions réelles (débit mesuré entre 1,3 et 4 tokens par seconde selon la charge concurrente). Le test comparatif avec Gemma-3-4B (section~\ref{sec:memoire}) offre une option de repli si cette latence devenait inacceptable pour les utilisateurs, mais n'a pas été retenue à ce stade sur décision explicite du porteur de projet.

\textbf{L'absence de superviseur de processus pour la couche applicative.} À plusieurs reprises durant le développement, le serveur \texttt{llama.cpp}, l'API et le serveur du widget se sont arrêtés entre deux sessions de travail (les conteneurs de stockage, gérés séparément avec redémarrage automatique, n'étant pas affectés), nécessitant un redémarrage manuel. La mise en place d'un superviseur de processus (par exemple \texttt{systemd}) pour la couche applicative est recommandée pour la mise en production définitive.

\textbf{Absence de dépôt de gestion de versions.} L'ensemble des modifications décrites dans ce rapport a été appliqué directement sur les fichiers du projet, sans historique de versions (\texttt{git}) -- un choix délibéré maintenu jusqu'ici, mais qui limite la traçabilité fine et la possibilité de revenir à un état antérieur précis. L'initialisation d'un dépôt est recommandée pour la suite du projet.
